\section*{ВСТУП}
\addcontentsline{toc}{section}{ВСТУП}

Сучасний розвиток програмних систем дедалі частіше пов'язаний з інтеграцією засобів штучного інтелекту, зокрема великих мовних моделей. Такі моделі використовуються для побудови діалогових асистентів, автоматизованого пояснення навчального матеріалу, генерації програмного коду, аналізу текстів та підтримки користувача під час роботи з інформаційними системами. Водночас практичне використання LLM потребує не лише самої моделі, а й прикладного програмного середовища, яке забезпечує введення запитів, збереження історії діалогів, відображення відповідей, керування параметрами генерації та обробку помилок.

Актуальність теми курсової роботи зумовлена необхідністю створення навчальної програмної системи, яка демонструє повний цикл взаємодії користувача з локальною великою мовною моделлю. На відміну від готових хмарних сервісів, власна реалізація дозволяє дослідити архітектуру такої системи, організацію web- та CLI-інтерфейсів, роботу зі сховищем даних, інтеграцію локальної pretrained-моделі, потокову генерацію відповідей та розширення через LoRA/QLoRA-адаптери.

Метою курсової роботи є розробка програмної системи діалогової взаємодії з великими мовними моделями з підтримкою web- та консольного інтерфейсів, збереженням історії діалогів у PostgreSQL та можливістю запуску локальних моделей через Hugging Face Transformers.

Для досягнення поставленої мети необхідно виконати такі завдання:
\begin{enumerate}
  \item проаналізувати предметну область діалогових систем на основі LLM;
  \item розглянути існуючі рішення та визначити їхні переваги й обмеження;
  \item сформувати функціональні та нефункціональні вимоги до системи;
  \item спроєктувати архітектуру програмного продукту та структуру бази даних;
  \item реалізувати web-інтерфейс для роботи з чатами та повідомленнями;
  \item реалізувати CLI-режим роботи;
  \item організувати збереження даних у PostgreSQL з використанням міграцій;
  \item інтегрувати локальний LLM-backend на основі Transformers;
  \item реалізувати runtime-перемикання моделей, streaming generation, зупинку генерації та підтримку LoRA/QLoRA-адаптерів;
  \item виконати тестування основних модулів системи.
\end{enumerate}

Об'єктом дослідження є процес побудови програмних систем для діалогової взаємодії користувача з великими мовними моделями.

Предметом дослідження є архітектурні, програмні та інтерфейсні засоби реалізації локальної LLM-системи з web- та CLI-режимами, базою даних PostgreSQL і локальним inference-backend.

Практичне значення роботи полягає у створенні працездатного програмного продукту, який може використовуватися як навчальний приклад побудови модульної LLM-системи, а також як основа для подальших експериментів з локальними моделями, адаптерами, параметрами генерації та розширенням функціональності.

\newpage
